% =========================================================
% Proof: Balls are open and closed
% Source: volume-iii/analysis/metric-spaces/notes/notes-balls.tex
% =========================================================

\subsection*{Balls are open and closed}
\label{prf:balls-open-closed}

\begin{remark}[Return]
\hyperref[prop:balls-open-closed]{$\leftarrow$ Back to Proposition (Balls are open and closed) in Notes}
\end{remark}

\begin{proposition}[Balls are open and closed]
Let $(X, d)$ be a metric space, $x \in X$, and $r > 0$.
\begin{enumerate}[label=(\roman*)]
  \item $B_r(x)$ is an open set.
  \item $\overline{B}_r(x)$ is a closed set.
\end{enumerate}
\end{proposition}

\begin{proof}
\Given A metric space $(X, d)$, a point $x \in X$, and a radius $r > 0$.
\Goal To show (i) $B_r(x)$ is open, and (ii) $\overline{B}_r(x)$ is closed.

\medskip
\noindent\textit{Part (i): $B_r(x)$ is open.}

\Strategy We show that every point of $B_r(x)$ has an open ball around it
still contained in $B_r(x)$.

Let $x_0 \in B_r(x)$.
By definition of open ball in a metric space,
\[
d(x,x_0) < r.
\]
Choose $\varepsilon>0$ with $\varepsilon < r - d(x,x_0)$. 
Claim $B_\varepsilon(x_0) \subseteq B_r(x)$.

Let $x_1 \in B_\varepsilon(x_0)$ be arbitrary. Then
\begin{align*}
d(x_0,x_1) &< \varepsilon,\\
d(x,x_1) &\le d(x,x_0) + d(x_0,x_1).
\end{align*}
Since $\varepsilon < r - d(x,x_0)$, we have $d(x,x_0) < r - \varepsilon$. Therefore,
\begin{align*}
d(x,x_1) 
&\le d(x,x_0) + d(x_0,x_1)\\
&< (r-\varepsilon) + \varepsilon\\
&= r.
\end{align*}
Hence, $x_1 \in B_r(x)$, and thus $B_\varepsilon(x_0) \subseteq B_r(x)$.

Since $x_1 \in B_\varepsilon(x_0)$ was arbitrary, it follows that
\[
B_\varepsilon(x_0) \subseteq B_r(x).
\]
Since $x_0 \in B_r(x)$ was arbitrary, every point of $B_r(x)$ admits an
$\varepsilon$-ball contained in $B_r(x)$. Therefore $B_r(x)$ is open.

\medskip
\noindent\textit{Part (ii): $\overline{B}_r(x)$ is closed.}

\Strategy We show that $X \setminus \overline{B}_r(x)$ is open, i.e.\ every
point outside the closed ball has an open ball around it disjoint from
$\overline{B}_r(x)$.

Let $x_0 \in X \setminus \overline{B}_r(x)$ be arbitrary. Then $d(x,x_0) > r$.

Choose $\varepsilon>0$ with $\varepsilon < d(x,x_0) - r$. 
Claim $B_\varepsilon(x_0) \subseteq X \setminus \overline{B}_r(x)$.

Let $x_1 \in B_\varepsilon(x_0)$ be arbitrary. Then
\begin{align*}
d(x_0,x_1) &< \varepsilon,\\
d(x,x_0) &\le d(x,x_1) + d(x_1,x_0).
\end{align*}

Since $\varepsilon < d(x,x_0) - r$, we have $\varepsilon + r < d(x,x_0)$. Therefore,
\begin{align*}
\varepsilon + r < d(x,x_0)
&\le d(x,x_1) + d(x_0,x_1)\\
&< d(x,x_1) + \varepsilon,
\end{align*}

which implies
\[
r < d(x,x_1).
\]

Hence $x_1 \notin \overline{B}_r(x)$, so
\[
B_\varepsilon(x_0) \subseteq X \setminus \overline{B}_r(x).
\]

Since $x_0 \in X \setminus \overline{B}_r(x)$ was arbitrary,
the complement of $\overline{B}_r(x)$ is open. Therefore
$\overline{B}_r(x)$ is closed.













\end{proof}

\begin{remark}[Proof shape]
% YOUR PROOF SHAPE HERE
\end{remark}

\begin{remark}[Dependencies]
\begin{itemize}
  \item Definition of open ball: \ref{def:open-ball}
  \item Definition of closed ball: \ref{def:closed-ball}
  \item Definition of open set: \ref{def:open-closed-set}
  \item Triangle inequality (metric axiom M4)
\end{itemize}
\end{remark}